% !PRINTED 252
\weberpage{252}
From
\[
-A_3\gamma_3 = x x_0 x_1 x_2\left(\frac{1}{x}+\frac{1}{x_0}+\frac{1}{x_1}+\frac{1}{x_2}\right)
\]
there results by (6) the constant value 1 for $A_3$. But in order to
determine $A_2$, we must go one term further in the development. In
\[
x_0^4 + \gamma_3 x_0 - 27 = -A_2 x_0^2
\]
we substitute for $x_0$ the development
\[
x_0 = -q^{-\frac{1}{3}} + 6\,q^{\frac{1}{3}} + \cdots,
\]
and find $A_2 = 18$, so that the equation sought for reads:
\weq{7}{x^4 + 18\,x^2 + \gamma_3\,x - 27 = 0.}

If $n$ is odd and not divisible by 3, let $c$ be assumed \gesperrt{divisible
by 12}.

It is then
\[
\lambda \equiv n, \quad \alpha \equiv 0, \quad \delta \equiv 0, \quad
\beta \equiv a\,\eth_2, \quad \gamma \equiv -\eth\,a_2 \pmod{12}
\]
and there results
\[
r^2 = (-1)^{\frac{\beta-1}{2}} = (-1)^{\frac{n-1}{2}}(-1)^{\frac{a-a'}{2}}, % [sic] printed a' here, not a_2 as in the parallel formulas above
\]
hence the theorem

4. \gesperrt{If $n$ is relatively prime to 6 and $c$ divisible by 12,
then the quantities}
\[
P^{2}_{c,\eth,a}\,\gamma_2(\omega)^{n-1}\,\gamma_3(\omega)^{\frac{n-1}{2}}
\]
\gesperrt{are roots of a multiplier equation of the first level.}

Since the functions $P^2$ become infinite or zero for no finite value
of $j(\omega)$, we conclude that the coefficients of the equation
whose roots are the $P^2$ are \gesperrt{polynomials in
$\gamma_2, \gamma_3$, and that the last coefficient is a constant.}

If $n$ is a prime $p$, the roots of this equation are
\[
x = p\left(\frac{\eta(p\omega)}{\eta(\omega)}\right)^2, \quad
x_h = (-1)^{\frac{p-1}{2}}\left(\frac{\eta\left(\dfrac{12h+\omega}{p}\right)}{\eta(\omega)}\right)^2
\quad h=0,1,\dots,p-1
\]
and the initial terms of the developments are the following:
\[
x = p\,q^{\frac{p-1}{6}}\dots, \quad
x_h = (-1)^{\frac{p-1}{2}}e^{\frac{2\pi i h}{p}}q^{-\frac{p-1}{6p}}\dots,
\]
whereby the value $(-1)^{\frac{p-1}{2}}p$ results for the last
coefficient. From the fact that the first term in the development of
%CONT>>
